\subsection*{Problems}

\textbf{13.5.} (a) Suppose X and Y are iid. random variables with finite mean. Prove that

E[X\vertX + Y]= (X + Y)/ 2. (Hint: Show that E [Y\vertX + Y]= E[X\vertX +Y ] .)

(b) Extend the result in part (a) to n iid. random variables.

\textbf{13.6.} Prove the conditional version of Jensen inequality: \phi(E [X\vert\mathcal{G}]) \leq E[\phi(X)\vert\mathcal{G}]

for convex function \phi(x) .

\textbf{13.7.} Suppose X and Y are random variables defined on the same probability space

and E [X2] < \infty .

(a) Prove the identity V ar(X) = Va r(E[X\vertY]) +E [(X - E[X\vertY])2].

(b) Show that the last term in part (a) can be interpreted as

E[(X- E [X\vertY])2]= E[E[X2\vertY]- E[X\vertY] 2].

\textbf{13.8.} (Generalized Bayes rule) Let (\Omega, \mathcal{G},P) be a probability space and \mathcal{G} be a

sub-\sigma-algebra of \mathcal{F} For any B \in \mathcal{G}and A \in \mathcal{F}with P(A) /=0, deduce that

\int

B P(A \vert\mathcal{G})dP\int

\Omega P(A \vert\mathcal{G})dP = P(B \capA)

P(A) .

Show that this equation reduces to the classical Bayes rule when \mathcal{G}is generated by

a partition of\Omega .

\textbf{13.9.} Suppose (X n)n\geq0 is a sequence of random variables defined on a probability

space (we take X 0 = 0 as the initial rv.)For n \geq 1, let

Yn \coloneqqX n -E [Xn\vertX0,X 1,...,X n- 1].

The random variable Y n can be interpreted as the new information contained in X n

relative to the information contained in the past X 0,...,X n- 1Let Sn denote the

partial sum S n \coloneqqY 1 +Y 2 +\cdot\cdot\cdot+ Yn for n \geq 1 and S 0 = 0. Show that (S n)n\geq0 is a

martingale relative to the filtration (\sigma(X 1,X 2,...,X n))n\geq0.

\textbf{13.10.} Let X n's be iid. random variables with finite mean \mu = E[X1]:

(a) Let T be an integer-valued random variable independent of (X n)n\geq1, and

assume T has finite expectation. Prove Wald's identity

E[X1 +X 2 +\cdot\cdot\cdot+ XT ]= \mu\cdot E [T] .

(b) Prove that Condition (iii) in Theorem 13.18 can be relaxed to E [T] < \infty and

E[\vertXn -X n- 1\vert\vert \mathcal{F}n]\leq c for all n .

(c) Let T be a stopping time for the filtration \mathcal{F}n \coloneqq\sigma(X 1,...,X n),f o r n \geq 1.

Apply part (b) to prove that Wald's identity holds if E [\vertX1\vert] < \infty and E [T] <

\infty .

\textbf{13.11.} (Polya's Urn Model)Consider an urn containing b black balls and w white

balls at time 0. At time i (for i \in N), we pick a ball uniformly at random from the

urn. If the color of the ball is black, we put two black balls back to the urn. If the

color is white, we put two white balls back. There are b +w+i balls in the urn after

step i Let Bi denote the number of black balls in the urn after step i :

(a) Compute the probabilities P(B 1 =b) and P(B 1 = b + 1).

(b) Let \mathcal{F}i be the \sigma -algebra generated by B 0,B 1,...,B i ,f o r i \geq 1. Let Z i =

Bi/(b + w+ i) be the fraction of black balls in the urn after the i -th step. Show

that (Z i)i\geq 0 is a martingale relative to the filtration (\mathcal{F}i)i\geq 0. (We define B 0 = b,

Z0 = b/(b + w), and \mathcal{F}0 ={ \emptyset,\Omega}.)

(c) Compute E [Zi] for i \geq 1.
